\part{A∞-\trCategory}
\markboth{I　A∞-\trCategory}{I　A∞-\trCategory}
本章叙述\trFukayaCategory{}底下的\trHomology{}代数。我们先对\trAinfty{}结构作一个快速
的回顾；这绝算不上是该理论的完整论述。大体上，唯一得到充分讨论的方面是\trIdentityMorphism，
因为它对\trFukayaCategory{}来说潜在地颇为棘手。此外，这套形式体系刻意保持朴实：这意味着
我们不使用 operad 或非交换几何，尽管这两种语言各自都能为这团公式乱麻带来不少启示。关于这些
材料有各种口味的替代文献，略举数条：一些奠基性的早期论文，比如 [70]、[62]，今天读来仍颇有
味道，尽管侧重之处与本书略有不同。较近的来源中，Keller 的工作 [82]、[83]、[84] 对基础材料给
出了平易近人的讲解，并有对代数与表示论问题的应用；Lefevre 的学位论文 [95] 则以\trHomotopy{}代
数的语言，从类似的视角给出了全面的论述。\trFukaya{}及其合作者写过好几篇专门面向\trSymplectic{}几
何的阐述 [51]、[52]、[53]、[58]、[60]、[56]、[55]；末三种文献的一般性程度远超我们的需要，
但愿意把这一点考虑在内的读者会发现，许多论证都可以缩减到我们这个更朴素的框架里。最后，
\trKontsevich{}的论文与讲义 [90]、[88]、[89] 中蕴含着大量洞见。

本章余下部分的总主题，是\trACategory{}中\trObject{}的构造与操作。我们将讨论\trMappingCone{}与%
\trExactTriangle，以及\trObject{}\trSplitting{}为\trDirectSummand{}的问题；还会考察一个更具体的构造，
即所谓一个\trObject{}沿另一\trObject{}的扭转。这一构造此前在文献中曾以两种不同的语境出现过。
其一是\trExceptionalCollection{}与\trMutation{}的概念，它作为在代数几何中构造刚性\trVectorBundle{}的一种手段而兴起：
见 [67] 与文集 [117]。该理论后来逐渐被推广到一般的\trTriangulated{}\trCategory（其各个方面
见 [65]、[26]、[27]、[28]）；一份有用的较新综述是 [66]。扭转构造的第二次、也是更晚近的出现，
则与球性\trObject{}有关——它在那里给出\trDerivedCategory{}的\trAutomorphism{} [134]。向%
\trACategory{}的推广由\trKontsevich{}开创，我们毫不客气地借用他的想法。除了搭建一般的形式体系，
我们还推导了若干推论，特别是\cref{thm:shifted-twist-split-generates}与\cref{thm:beilinson-spectral-sequence}，它们在日后与\trSymplectic{}几何及%
\trLefschetzFibration{}建立联系时会派上用场。本章的收尾略显技术性：简短讨论了带 $\ZZ/2$ 对
称性的\trACategory。
